% ============================================================
% CONCLUSION
% ============================================================

\section*{Results and Conclusion}

% ----------------------------------------------------------
% Classical Thickness Distributions
% ----------------------------------------------------------

\subsection*{Excel Model Capabilities}

The developed Excel-based simulation tool provides a parametric visualization of NACA airfoil geometries. 
The model allows real-time generation of the airfoil profile based on user-defined input parameters within a designated control panel.

The tool is capable of:
\begin{itemize}
\item Displaying the complete airfoil geometry (the upper and lower surfaces).
\item Plotting the mean camber line for cambered airfoils.
\item Indicating the leading-edge radius based on the selected series formulation.
\item Updating the airfoil profile dynamically as geometric parameters are modified.
\end{itemize}

The geometry responds directly to variations in defining parameters 
such as maximum camber, camber position, and thickness ratio. 
Any change in the control area automatically updates the plotted airfoil, 
enabling rapid comparative analysis between different configurations.

This parametric structure allows efficient investigation of geometric trends and facilitates understanding of how individual parameters influence airfoil shape. 
Although the tool focuses primarily on geometric modeling rather than full aerodynamic performance prediction, 
it provides a clear and practical platform for preliminary airfoil design and educational analysis.

% ----------------------------------------------------------
% Classical Thickness Distributions
% ----------------------------------------------------------

\subsection*{Comparison and Assessment}

This study summarizes the evolution of NACA airfoils, 
from early geometry-based definitions to modern aerodynamic and empirical design approaches. 
Initially, airfoil design relied primarily on geometric parameters 
such as camber, camber position, and thickness, which allowed engineers to generate airfoil shapes using analytical equations. 
These early methods provided a foundation for understanding airfoil geometry and its influence on aerodynamic performance.

However, 
with the advancement of aerodynamic theory and computational tools, airfoil design has gradually shifted toward performance-based approaches. 
Modern airfoil development focuses on achieving specific aerodynamic objectives, 
such as target lift coefficients, favorable pressure distributions, and improved lift-to-drag ratios.

Simulation also plays a critical role in airfoil analysis, 
although numerical results may contain errors due to discretization, interpolation, numerical approximation, and modeling assumptions. 
Therefore, simulation results must be carefully validated and interpreted.

In practical applications, 
the digits in the NACA series are not strictly followed during manufacturing. Although airfoils may retain their NACA designation, 
their geometries are often refined through aerodynamic analysis, numerical simulations, and experimental validation. 
These refinements ensure that the final airfoil meets performance, stability, and efficiency requirements for specific applications such as aircraft wings, propellers, or turbine blades.
